%% GENERATED by tools/make_aggregation_robustness.py -- do not hand-edit.

\section{Does the headline depend on how the ratio was averaged?}
\label{app:aggregation}

A ratio is a nonlinear function of two means, so its value depends on the order
of averaging: $\E[X/Y] \ge \E[X]/\E[Y]$ by Jensen, with a gap that grows
with the spread. Table~\ref{tab:pointwise} reports the mean of per-cell
ratios, aggregated per seed. That is the right inferential unit --- the twelve
noise levels share a training set and a fitted model --- but being the right
unit does not make it the only defensible summary, so here is the same data
under five.

\begin{center}\small
\captionof{table}{The headline under five estimands. (1) is what
Table~\ref{tab:pointwise} reports. The interval is a paired bootstrap over
seeds, 20,000 resamples; seeds are resampled rather than cells, because cells
within a seed share a fitted model.}
\label{tab:aggregation}
\begin{tabular}{rccccc}
\toprule
$\nseq$ & (1) per-cell & (2) ratio of means & (3) geometric & (4) median
        & (5) paired 95\% CI \\
\midrule
32 & $7.0$ & $6.8$ & $6.1$ & $6.8$ & $[6.1,\,8.1]$ \\
64 & $7.8$ & $6.6$ & $6.7$ & $6.6$ & $[6.3,\,9.5]$ \\
128 & $8.8$ & $8.0$ & $7.8$ & $8.6$ & $[7.9,\,9.6]$ \\
256 & $9.7$ & $8.8$ & $8.6$ & $9.2$ & $[8.5,\,10.9]$ \\
512 & $9.0$ & $8.6$ & $7.9$ & $8.2$ & $[8.1,\,9.9]$ \\
1024 & $12.5$ & $9.7$ & $10.8$ & $11.1$ & $[10.5,\,14.7]$ \\
2048 & $15.5$ & $10.0$ & $13.2$ & $12.9$ & $[13.0,\,18.9]$ \\
4096 & $17.6$ & $12.4$ & $15.6$ & $16.0$ & $[15.4,\,19.9]$ \\
8192 & $20.2$ & $16.1$ & $18.8$ & $19.5$ & $[18.3,\,22.0]$ \\
\bottomrule
\end{tabular}
\end{center}

\noindent Two things follow, and the second is the one worth stating plainly.

The conclusion does not depend on the choice: every estimand at every size
exceeds $6.1$, every bootstrap interval sits well clear of $1$, and no
aggregation produces a reversal anywhere.

But \textbf{the estimand we report is systematically the largest of the
four}, and by a margin that grows with $\nseq$ --- up to $5.5$ at
$\nseq=2048$. The ratio of per-seed averaged errors, which is the summary a
reader is most likely to have in mind, runs $6.6$ to $16.1$
rather than $\ratiolo$ to $\ratiohi$. This is Jensen's inequality doing
exactly what it must, not an error, but quoting only the largest of five
summaries would be a choice presented as a fact.

Resolving the average over noise levels shows where the aggregation hides
structure: the ratio is not flat in $t$ but peaks near $t \approx 0.5$ and
falls at both ends, and the weakest single (size, level) cell in the whole
grid is $3.6$. The estimator's advantage is largest at moderate
noise, where the posterior is neither dominated by the likelihood nor by the
prior --- which is where knowing the transition structure should help most.
